% A proof of Graffiti.pc Conjecture 141 (WOWII): t(G) >= floor(g/2) - 1 + max_v l(v),
% via the stronger bound t >= Delta + g - 3 for triangle-free cyclic graphs.
\documentclass[11pt]{article}
\ifdefined\XeTeXversion\else\pdfoutput=1\fi
\usepackage[margin=1.1in]{geometry}
\usepackage{amsmath,amssymb,amsthm}
\usepackage[colorlinks=true,linkcolor=blue,citecolor=blue,urlcolor=blue]{hyperref}
\Urlmuskip=0mu plus 1mu\relax
\setlength{\emergencystretch}{3em}
\numberwithin{equation}{section}

\newtheorem{theorem}{Theorem}[section]
\newtheorem{lemma}[theorem]{Lemma}
\newtheorem{corollary}[theorem]{Corollary}
\newtheorem{proposition}[theorem]{Proposition}
\theoremstyle{remark}
\newtheorem{remark}[theorem]{Remark}

\newcommand{\tr}{\operatorname{t}}
\newcommand{\gr}{\operatorname{g}}
\newcommand{\lv}{\ell}

\title{Largest induced trees in graphs of given girth and maximum degree,\\
with a proof of Graffiti.pc's Conjecture 141}
\author{Alper Ferudun\\\texttt{alper@mercurycodelab.com}}
\date{July 2026}

\begin{document}
\maketitle

\begin{abstract}
For a finite simple graph $G$, let $\tr(G)$ denote the largest number of
vertices inducing a tree, $\gr(G)$ the girth, $\Delta(G)$ the maximum degree,
and for a vertex $v$ let $\lv(v)$ denote the independence number of the
subgraph induced by the neighbourhood of $v$.  We prove that every finite
connected triangle-free graph containing a cycle satisfies
$\tr(G)\ge\Delta(G)+\gr(G)-3$, with equality for all cycles $C_g$ with
$g\ge4$ and all complete bipartite graphs $K_{a,b}$ with $a,b\ge2$.  As a corollary we prove Conjecture~141 of the
conjecture-making program Graffiti.pc (DeLaVi\~na, \emph{Written on the Wall
II}, 2005): every finite connected graph satisfies
$\tr(G)\ge\lfloor\gr(G)/2\rfloor-1+\max_v\lv(v)$.  The proofs are elementary
maximality arguments.  Both results have been formalized and machine-checked
in Lean~4 against the statement of the conjecture in the Google DeepMind
\emph{Formal Conjectures} repository.
\end{abstract}

\section{Introduction}

Graffiti.pc, developed by DeLaVi\~na~\cite{DeL01,DeLhistory} as a successor of
Fajtlowicz's program Graffiti~\cite{Faj88}, generates conjectural inequalities
between graph invariants; its conjectures are collected in the lists
\emph{Written on the Wall II}~\cite{WOWII}, with a curated register maintained
by West~\cite{WestReg}.  A run of these conjectures bounds from below the
invariant $\tr(G)$, the largest order of an induced subtree, whose systematic
study goes back to Erd\H{o}s, Saks, and S\'os~\cite{ESS86}.

Throughout, $G$ is a finite simple graph, $\gr(G)$ its girth, $\Delta(G)$ its
maximum degree, and for a vertex $v$ we write
$\lv(v)$ for the independence number of the subgraph induced by the
neighbourhood $N(v)$.  Conjecture~141 of Written on the Wall II, dated 2005,
asserts~\cite{WOWII}:

\begin{quote}
\emph{If $G$ is a simple connected graph, then
$\tr(G)\ \ge\ \lfloor \gr(G)/2\rfloor - 1 + \max_v \lv(v)$.}
\end{quote}

At the time of writing the conjecture is listed as open on DeLaVi\~na's
register (status ``O'')~\cite{WOWII} and is stated as a research-open item in
the Google DeepMind \emph{Formal Conjectures} repository~\cite{FC,FC141}.  We
prove it.
The main step is the following sharper bound, which does not involve
$\lv$ at all and may be of independent interest.

\begin{theorem}\label{thm:main}
Let $G$ be a finite simple connected triangle-free graph that contains a
cycle.  Then
\[
  \tr(G)\;\ge\;\Delta(G)+\gr(G)-3 .
\]
Equality holds for every cycle $C_g$ ($g\ge4$) and every complete bipartite
graph $K_{a,b}$ with $a,b\ge2$.
\end{theorem}

\begin{corollary}[Conjecture 141]\label{cor:141}
Every finite simple connected graph $G$ on at least two vertices satisfies
$\tr(G)\ \ge\ \lfloor \gr(G)/2\rfloor - 1 + \max_v \lv(v)$, where
$\gr(G):=0$ if $G$ is acyclic.
\end{corollary}

The deduction of the corollary is short: for girth $\le5$ (and for acyclic
graphs) an induced star at a vertex maximizing $\lv$ already beats the bound,
while for girth $\ge6$ the graph is triangle-free, so $\lv(v)=\deg(v)$ for
every $v$, and Theorem~\ref{thm:main} applies with
$\gr-3\ge\lfloor\gr/2\rfloor-1$.

The proof of Theorem~\ref{thm:main} is a maximality argument in the spirit of
our proof of Graffiti.pc's Conjecture~143~\cite{W143}: a maximum-order induced
tree containing the closed neighbourhood of a maximum-degree vertex either
spans (impossible in a cyclic graph) or admits an outside vertex with two
neighbours on the tree; the tree path joining those two neighbours closes a
cycle, hence has length at least $\gr(G)-2$, and it can share at most three
vertices with the closed neighbourhood.

\paragraph{Related work.}
The foundational paper~\cite{ESS86} bounds $\tr(G)$ in terms of order, size,
radius, and independence and clique numbers; we are not aware of a published
lower bound for $\tr$ combining girth with the maximum degree or with
neighbourhood independence, and the closest Graffiti.pc antecedents (the
induced-forest bounds of DeLaVi\~na--Waller~\cite{DW04}; see also
Hertz--Marcotte--Schindl~\cite{HMS14}) concern induced forests rather than
trees.  DeLaVi\~na's resolved-conjecture lists do not contain
Conjecture~141~\cite{WOWII}.  As with any short elementary argument, we do
not claim it could not have been observed before.

\paragraph{Verification.}
Theorem~\ref{thm:main} and Corollary~\ref{cor:141} have been formalized and
machine-checked in Lean~4~\cite{Lean4} on top of Mathlib~\cite{mathlib},
against the pre-existing formal statement \texttt{conjecture141} of the
Formal Conjectures repository~\cite{FC141}; see Section~\ref{sec:formal}.
Independently, the conjecture and all intermediate claims were verified
exhaustively by exact computation on all $995$ connected graphs on $2$--$7$
vertices and on structured families (cages, prisms, generalized theta graphs,
cycles with pendant trees, complete bipartite graphs, and seeded random
graphs); no violation was found, and every equality case located has girth
exactly~$4$.  The computations play no role in the proofs.

\section{Notation}

All graphs are finite and simple.  For $S\subseteq V(G)$, $G[S]$ is the
induced subgraph; $S$ \emph{induces a tree} if $G[S]$ is connected and
acyclic, and $\tr(G)=\max\{|S|:G[S]\text{ is a tree}\}$.  The girth of a
graph containing a cycle is the minimum length of a cycle; following the
Mathlib convention used by the formal statement, an acyclic graph has girth
$0$.  $N(v)$ and $N[v]=N(v)\cup\{v\}$ are the open and closed
neighbourhoods; $\lv(v)$ is the independence number of $G[N(v)]$.  Note
$\lv(v)\le\deg(v)$ always, with equality when $N(v)$ is independent, i.e.\
for all $v$ when $G$ is triangle-free.

\section{The star lemma}

\begin{lemma}\label{lem:star}
For every vertex $v$ of any graph $G$: $\tr(G)\ \ge\ \lv(v)+1$.
\end{lemma}

\begin{proof}
Let $S\subseteq N(v)$ be independent with $|S|=\lv(v)$.  Then
$G[\{v\}\cup S]$ is a star: all edges $vs$ ($s\in S$) are present and there
are no edges inside $S$.  A star is a tree.
\end{proof}

\section{Proof of Theorem~\texorpdfstring{\ref{thm:main}}{1.1}}

\begin{proof}[Proof of Theorem~\ref{thm:main}]
Write $g=\gr(G)\ge4$ and $\Delta=\Delta(G)$.  Fix a vertex $v$ of degree
$\Delta$.  Since $G$ is triangle-free, $N(v)$ is independent, so by the proof
of Lemma~\ref{lem:star} the closed neighbourhood $N[v]$ induces a star, hence
a tree.  Among all sets $S\supseteq N[v]$ inducing a tree choose one, say
$S^*$, of maximum cardinality, and let $T=G[S^*]$.

$S^*$ is a proper subset of $V(G)$: an induced spanning tree would make $G$
itself a tree, contradicting the existence of a cycle.  Since $G$ is
connected, some vertex $z\notin S^*$ has a neighbour in $S^*$; and $z$ must
have two distinct neighbours $a,b\in S^*$, for otherwise $G[S^*\cup\{z\}]$
would be the tree $T$ with a pendant vertex attached, an induced tree
containing $N[v]$ and larger than $T$ --- contradicting maximality.

Let $P$ be the unique $a$--$b$ path in $T$.  The edges of $P$ together with
$za$ and $zb$ form a cycle of length $|E(P)|+2$, so
\begin{equation}\label{eq:girthbound}
  |E(P)|\;\ge\;g-2 .
\end{equation}

Next, $|V(P)\cap N[v]|\le 3$.  Every edge of $G$ between two vertices of
$S^*$ is an edge of the tree $T$ (as $T$ is induced), and paths between fixed
endpoints in a tree are unique.  If $v\notin V(P)$ and $P$ contained two
distinct vertices $u_1,u_2\in N(v)$, then the segment of $P$ between $u_1$
and $u_2$ would be the unique $u_1$--$u_2$ path in $T$; but $u_1vu_2$ is also
a $u_1$--$u_2$ path in $T$, so the segment would pass through $v$,
contradicting $v\notin V(P)$; hence $|V(P)\cap N[v]|\le1$ in this case.  If
$v\in V(P)$, then for any $u\in V(P)\cap N(v)$ the segment of $P$ between $v$
and $u$ is the unique $v$--$u$ path of $T$, which is the single edge $vu$;
so $u$ is adjacent to $v$ \emph{along} $P$, and a path has at most two edges
at any vertex: $|V(P)\cap N(v)|\le2$, so $|V(P)\cap N[v]|\le3$.

Finally, count.  $V(P)$ and $N[v]$ are subsets of $S^*$, so by
inclusion--exclusion and~\eqref{eq:girthbound},
\[
  |S^*|\;\ge\;|V(P)\cup N[v]|\;=\;|V(P)|+|N[v]|-|V(P)\cap N[v]|
  \;\ge\;(g-1)+(\Delta+1)-3\;=\;\Delta+g-3 .
\]
Since $T$ is an induced tree, $\tr(G)\ge|S^*|\ge\Delta+g-3$.

For sharpness: in $C_g$ ($\Delta=2$) the largest induced trees are the paths
obtained by deleting one vertex, so $\tr=g-1=\Delta+g-3$; in $K_{a,b}$ with
$a\le b$ ($g=4$, $\Delta=b$) every induced tree is a star (two vertices on
each side already induce $C_4$), so $\tr=b+1=\Delta+g-3$.
\end{proof}

\begin{remark}
Under the hypotheses of Theorem~\ref{thm:main}, an alternative proof of the
slightly weaker bound $\tr\ge\Delta+\lfloor g/2\rfloor-1$ (still sufficient
for Corollary~\ref{cor:141}) is worth recording.  For
$r=\lfloor g/2\rfloor-1$, the ball $B(v,r)$ induces a tree in any graph of
girth $g$: every $G$-edge inside the ball is an edge of a breadth-first
search tree rooted at $v$, since a non-tree edge would close a cycle of
length at most $2r+1<g$ through the last common ancestor of its endpoints.
Now take $v$ of maximum degree; triangle-freeness gives $g\ge4$, so
$r\ge1$.  If $G$ contains a cycle the ball is proper, so all $r$ distance
layers are nonempty and $|B(v,r)|\ge1+\Delta+(r-1)$.
\end{remark}

\section{Proof of Corollary~\texorpdfstring{\ref{cor:141}}{1.2}}

\begin{proof}[Proof of Corollary~\ref{cor:141}]
Let $L=\max_v\lv(v)$; note $L\ge1$ since $G$ is connected on $\ge2$
vertices.  If $G$ is acyclic then $\gr=0$ and the claimed bound is
$L-1\le\tr$, which follows from Lemma~\ref{lem:star}.  If
$3\le\gr\le5$ then $\lfloor \gr/2\rfloor-1\le1$ and Lemma~\ref{lem:star}
again gives $\tr\ge L+1\ge\lfloor \gr/2\rfloor-1+L$.  If $\gr\ge6$ then $G$
is triangle-free, so $\lv(v)=\deg(v)$ for every $v$ and $L=\Delta$; by
Theorem~\ref{thm:main},
$\tr\ge\Delta+\gr-3\ge\Delta+\lfloor \gr/2\rfloor-1=L+\lfloor
\gr/2\rfloor-1$, using $\gr-3\ge\lfloor \gr/2\rfloor-1$ for $\gr\ge4$.
\end{proof}

\begin{remark}
The corollary is sharp precisely in girth~$4$.  Equality holds for $C_4$, for
every $K_{a,b}$ ($a,b\ge2$), and for the infinite family obtained from $C_4$
by attaching $k\ge0$ pendant vertices to one vertex ($\lv_{\max}=k+2$,
$\tr=k+3$); an exhaustive check of all connected graphs on at most seven
vertices finds no equality case of any other girth.  That equality
\emph{requires} girth~$4$ follows from the results above: for acyclic $G$
and for $g=3$ the star bound $\tr\ge L+1$ beats the right-hand side by at
least $2$ resp.\ $1$; for $g=5$ triangle-freeness gives $L=\Delta$ and
Theorem~\ref{thm:main} yields $\tr\ge\Delta+2=L+\lfloor5/2\rfloor$, one more
than required; and for $g\ge6$ the inequality $g-3>\lfloor g/2\rfloor-1$ is
strict.
\end{remark}

\section{Formalization}\label{sec:formal}

The Formal Conjectures project~\cite{FC} maintains Lean~4 formalizations of
open conjectures.  Conjecture~141 appears in \texttt{GraphConjecture141.lean}
as \texttt{conjecture141}, marked \texttt{research open}~\cite{FC141}, in
the denominator-explicit integer form
\[
  \lfloor\gr(G)/2\rfloor-1+\max_v\lv(v)\le\tr(G).
\]
Here $\tr$ and $\lv$ are the repository definitions
\texttt{largestInducedTree\allowbreak Size} and
\texttt{indepNeighbors\allowbreak Card}; Mathlib girth is natural-valued and
equals zero on acyclic graphs.

We have produced a complete, \texttt{sorry}-free Lean~4 proof of exactly this
statement, machine-checked with Lean toolchain v4.27.0 against current
Mathlib, together with the supporting API: the star construction, the
maximum-induced-tree selection with prescribed vertices, the two-neighbour
maximality obstruction, the tree-path girth certificate, and the
three-vertex overlap bound of Theorem~\ref{thm:main}.  The development
builds on the reusable induced-tree API we contributed alongside
Conjecture~143~\cite{W143}.  The proof uses no \texttt{native\_decide} and no
additional axioms (\texttt{\#print axioms} reports \texttt{propext},
\texttt{Classical.choice}, \texttt{Quot.sound}).  The source is available as
pull request \#4454 to the Formal Conjectures repository (which also contains
our formalization of Conjecture~143) and as ancillary files with this
submission.

\paragraph{Acknowledgements.}
The author thanks the maintainers of the Formal Conjectures repository for
the formal statements, and Ermelinda DeLaVi\~na and Douglas B.\ West for
maintaining the Graffiti.pc conjecture lists.

\begin{thebibliography}{10}

\bibitem{DeL01}
E.~DeLaVi\~na,
\emph{Graffiti.pc},
Graph Theory Notes of New York XLII:3 (2002), 26--30.

\bibitem{DeLhistory}
E.~DeLaVi\~na,
\emph{Some history of the development of Graffiti},
in: Graphs and Discovery, DIMACS Ser.\ Discrete Math.\ Theoret.\ Comput.\
Sci.\ 69, AMS, 2005, 81--118.

\bibitem{WOWII}
E.~DeLaVi\~na,
\emph{Written on the Wall II: Conjectures of Graffiti.pc},
web list, University of Houston--Downtown,
\url{http://cms.dt.uh.edu/faculty/delavinae/research/wowII/} (retrieved July
2026).

\bibitem{WestReg}
D.~B. West,
\emph{Some conjectures of Graffiti.pc},
web register,
\url{https://dwest.web.illinois.edu/regs/graffiti.html} (retrieved July
2026).

\bibitem{Faj88}
S.~Fajtlowicz,
\emph{On conjectures of Graffiti},
Discrete Math.\ 72 (1988), 113--118.

\bibitem{ESS86}
P.~Erd\H{o}s, M.~Saks, and V.~T. S\'os,
\emph{Maximum induced trees in graphs},
J.~Combin.\ Theory Ser.\ B 41 (1986), 61--79.

\bibitem{DW04}
E.~DeLaVi\~na and B.~Waller,
\emph{On some conjectures of Graffiti.pc on the maximum order of induced
subgraphs},
Congr.\ Numer.\ 166 (2004), 11--32.

\bibitem{HMS14}
A.~Hertz, O.~Marcotte, and D.~Schindl,
\emph{On the maximum orders of an induced forest, an induced tree, and a
stable set},
Yugosl.\ J.\ Oper.\ Res.\ 24 (2014), no.~2, 199--215.
\href{https://doi.org/10.2298/YJOR130402037H}{doi:10.2298/YJOR130402037H}.

\bibitem{W143}
A.~Ferudun,
\emph{Largest induced trees, girth, and the second-smallest degree: a proof
of Graffiti.pc's Conjecture 143},
Lean~4 formalization, pull request \#4454 to the Google DeepMind
\emph{Formal Conjectures} repository,
\url{https://github.com/google-deepmind/formal-conjectures/pull/4454}
(2026).

\bibitem{FC}
Google DeepMind,
\emph{Formal Conjectures},
\href{https://github.com/google-deepmind/formal-conjectures}{GitHub repository}
(retrieved July 2026).

\bibitem{FC141}
Formal Conjectures,
\emph{GraphConjecture141.lean},
\href{https://github.com/google-deepmind/formal-conjectures/blob/main/FormalConjectures/WrittenOnTheWallII/GraphConjecture141.lean}{source file}
(retrieved July 2026).

\bibitem{Lean4}
L.~de~Moura and S.~Ullrich,
\emph{The Lean 4 theorem prover and programming language},
in: Automated Deduction -- CADE 28, LNCS 12699, Springer, 2021, 625--635.

\bibitem{mathlib}
The mathlib Community,
\emph{The Lean mathematical library},
in: Proc.\ 9th ACM SIGPLAN Int.\ Conf.\ Certified Programs and Proofs
(CPP 2020), ACM, 2020, 367--381.

\end{thebibliography}

\end{document}
